% =====================================================================
%  Цель и задание лабораторной работы.
% =====================================================================

\textbf{Цель:} изучить и отработать практические навыки применения
методов автоматического реферирования текстовых документов.

\vspace{1em}

\textbf{Задание (вариант 16)}: разработать систему автоматического
реферирования текстовых документов со следующей комбинацией признаков
варианта:
\begin{itemize}
  \item \textbf{Вид реферата} -- классический реферат (выборка наиболее
  значимых предложений документа) и реферат в виде иерархического
  списка ключевых слов и словосочетаний;
  \item \textbf{Реализуемые методы отбора предложений} -- алгоритм из
  методички (модифицированный TF-IDF с учётом положения предложения в
  документе и абзаце), TextRank и метод на эмбеддингах предложений;
  \item \textbf{Улучшение реферата} -- перефразирование отобранных
  предложений с помощью LLM без добавления новых фактов.
\end{itemize}

\vspace{1em}

Согласно методическим указаниям, для каждого метода отбора предложений
необходимо:
\begin{enumerate}
  \item рассчитать вес каждого предложения документа в соответствии с
  правилом метода;
  \item отобрать заданное число предложений с наибольшим весом,
  сохранив их исходный порядок следования в документе;
  \item отдельно построить реферат-список ключевых слов: термины
  документа с наибольшим весом, при возможности сгруппированные в
  иерархию (корневое слово -- словосочетания, в которые оно входит),
  например: лазер -- красный лазер, синий лазер;
  \item сравнить результаты разных методов по числу отобранных
  предложений, степени сжатия и затраченному времени.
\end{enumerate}

Кроме того, система должна:
\begin{itemize}
  \item на выходе выдавать активную ссылку на документ, текст реферата
  (классического и в виде ключевых слов) и сводную статистику
  сравнения методов по всей тестовой коллекции;
  \item предоставлять средства сохранения результата в файл и печати;
  \item иметь простой и доступный пользовательский интерфейс со
  справочными подсказками.
\end{itemize}
